% Source: https://artofproblemsolving.com/community/c3h1484950_triangles_problem
% Stable identity: aops-topic-1484950-q01
Let $\angle C=x$.

\begin{asy}
usepackage("amsmath");
size(6cm);
defaultpen(fontsize(10pt));
real c=28;
pair C=(0,0), D=(1,0), Ee=2*Cos(c)*dir(c), B=(1+2*Cos(2c),0), A=2.9*dir(c);
draw(A--B--C--cycle,linewidth(0.9));
draw(B--Ee--D,blue);
path tk(pair P,pair Q){pair m=(P+Q)/2; pair u=0.05*unit(rotate(90)*(Q-P)); return (m-u)--(m+u);}
draw(tk(B,Ee),blue); draw(tk(Ee,D),blue); draw(tk(D,C),blue);
real a1=degrees(B-Ee); if(a1>180) a1-=360; real a2=degrees(A-Ee); if(a2>180) a2-=360;
draw(arc(Ee,0.22,a1,a2),red);
label("$84^\circ$",Ee+0.4*dir((a1+a2)/2),red);
draw(arc(C,0.25,0,c),gray); label("$x$",C+0.37*dir(c/2),gray);
dot(A);dot(B);dot(C);dot(D);dot(Ee);
label("$A$",A,N); label("$B$",B,SE); label("$C$",C,SW); label("$D$",D,S); label("$E$",Ee,dir(118));
label("\mbox{}",(min(currentpicture,user=true)+max(currentpicture,user=true))/2,black); // keeps the text after the figure black when LaTeX typesets the labels inline
\end{asy}

\begin{itemize}
\item Triangle $DEC$ is isosceles with $DE=DC$, so $\angle DEC=\angle DCE=x$. The exterior angle at $D$ is $\angle EDB=\angle DEC+\angle DCE=2x$.
\item Triangle $BED$ is isosceles with $EB=ED$, so $\angle EBD=\angle EDB=2x$.
\item In triangle $EBC$, the angle $\angle AEB$ is the exterior angle at $E$, so $\angle AEB=\angle EBC+\angle ECB=2x+x=3x$.
\end{itemize}
So $3x=84^\circ$ and $\angle C=\boxed{28}$ degrees.

This is why the book's solution labels the angles $x$, $2x$, $2x$: no new theorem is involved, only the base angles of isosceles triangles and the exterior angle theorem, used three times.
